\bookchapter{Глава 39. Инженерия контекста — помогите ИИ понять ваши намерения}

\chapterintro{Знакомо ли вам это: открыли новое окно диалога, и ИИ будто амнезию получил — забыл всё о вашем проекте? Это не ИИ вас нарочно дразнит, просто ему от природы свойственна «забывчивость». Эта глава научит ставить ИИ «систему долговременной памяти» и заодно познакомит с одной из самых горячих парадигм разработки 2026 года — разработкой по спецификациям.}

\booksection{Файл памяти проекта: «долговременная память» ИИ}

У ИИ-моделей есть фундаментальное ограничение: их память «краткосрочна». В начале каждого диалога ИИ знает о вашем проекте ровно ноль. Вам приходится раз за разом сообщать ему «у нас React», «бэкенд на Node.js», «база данных PostgreSQL», «строки в коде берём в двойные кавычки». Это не только неэффективно, но и чревато упущениями — как заново проводить онбординг новому коллеге каждое утро.

Решение — «файл памяти проекта»: текстовый файл в корне проекта, который ИИ-инструменты автоматически читают в начале каждого диалога. У разных инструментов свои соглашения об именах:

\begin{center}
\begin{tabularx}{\linewidth}{|X|X|}
\hline
\textbf{Инструмент} & \textbf{Имя файла} \tabularnewline
\hline
Cursor & \texttt{.\allowbreak cursorrules} \tabularnewline
Claude Code & \texttt{CLAUDE.\allowbreak md} \tabularnewline
GitHub Copilot & \texttt{copilot-\allowbreak instructions.\allowbreak md} \tabularnewline
Универсальный вариант & \texttt{AGENTS.\allowbreak md} \tabularnewline
\hline
\end{tabularx}
\end{center}

\begin{center}\small\color{gray}Соглашения об именах файлов памяти проекта в разных инструментах\end{center}

Имена разные, но принцип один: записывайте в этот файл то, что ИИ должен «помнить всегда». Хороший файл памяти проекта должен содержать: обзор проекта (что это за проект), техстек (какие технологии используются), стандарты кода (имена, отступы, стиль комментариев), структуру каталогов (как организован код), соглашения разработки (стратегия веток, правила коммитов, требования к тестам). CLAUDE.md у Karpathy — всего 70 строк, но покрывает всю ключевую информацию: лаконично и полно.

\begin{lstlisting}
# Проект: личное приложение для учёта финансов
## Обзор
Полностековое приложение для личных финансов: помогает записывать ежедневные доходы и расходы.
## Технологии
- Фронтенд: React + TypeScript + Tailwind CSS
- Бэкенд: Node.js + Express + TypeScript
- База данных: SQLite (разработка) / PostgreSQL (прод)
- Тесты: Vitest
## Стиль кода
- Используем двойные кавычки
- Отступ в 2 пробела
- Имена функций в camelCase
- Типы в PascalCase
- У каждой функции — комментарий JSDoc
## Структура каталогов
/src
  /frontend  - код фронтенда
  /backend   - код бэкенда
  /shared    - общие типы фронта и бэка
  /tests     - файлы тестов
## Правила разработки
- Сообщения коммитов на русском
- Для новых функций сначала пишем тесты
- API всегда возвращает { success, data, error }
\end{lstlisting}

\begin{tcbwarn}{Внимание: не превращайте его во «второй документ требований»}
Следите за длиной этого файла — не пишите из него многотомный документ требований на десятки страниц. Задача файла памяти — «дать ИИ быстрый фон», а не «детально описать каждую функцию». Держите его кратким, структурированным, только то, «что ИИ должен знать».

\end{tcbwarn}

\booksection{Разработка по спецификациям (SDD)}

Spec-Driven Development (SDD, разработка по спецификациям) — парадигма разработки, взрослеющая с 2026 года. Её суть в одном предложении: сначала спецификация, потом код. В отличие от классического TDD (разработки через тестирование), «спецификация» в SDD — это не только тест-кейсы, но и описание требований, определения интерфейсов, модели данных и т. д. ИИ генерирует код по этим спецификациям, а обязанность человека — писать и поддерживать спецификации.

Рабочий процесс SDD состоит из пяти шагов:

\begin{enumerate}
  \item Опишите требования на естественном языке (в файле .md);
  \item Определите интерфейсы и структуры данных (в файле спецификации);
  \item Поручите ИИ сгенерировать реализацию кода по спецификации;
  \item Запустите тесты, чтобы проверить соответствие результата спецификации;
  \item Если не соответствует — правьте спецификацию или пусть ИИ правит код.
\end{enumerate}

В этом процессе файл спецификации — «единственный источник истины»: код можно перегенерировать, но спецификация сначала должна быть верной.

Написать файл спецификации

Достаточна ли спецификация?

Достаточна

ИИ генерирует код по спецификации

Прогнать тесты и проверить результат

Недостаточна / тесты не проходят

Проблема в коде правим код; проблема в спецификации правим спецификацию

Снова к проверке

Всё пройдено сдаём

\begin{center}\small\color{gray}Рисунок 3: рабочий процесс разработки по спецификациям (SDD)\end{center}

Сила SDD в том, что файл спецификации стабильнее кода. Когда меняется модель или инструмент, ваш файл спецификации остаётся в силе — достаточно поручить новому ИИ-инструменту перегенерировать код. Именно поэтому книга подчёркивает «методология важнее инструментов»: файл спецификации — конкретный носитель вашей методологии.

\booksection{Управление окном контекста}

Даже с файлом памяти проекта вы всё равно столкнётесь в диалогах с проблемой «контекста не хватает». Окно контекста ИИ-модели (объём информации, обрабатываемой одновременно) ограничено. Когда диалог слишком длинный или файлов упомянуто слишком много, ИИ «забывает» раннюю информацию, и качество результата падает — как у друга, с которым вы проговорили шесть часов и который давно не помнит, о чём вы сказали в самом начале.

Ключевая стратегия управления окном контекста — «послойные ссылки»:

\begin{center}
\begin{tabularx}{\linewidth}{|X|X|X|}
\hline
\textbf{Слой контекста} & \textbf{Содержимое} & \textbf{Способ управления} \tabularnewline
\hline
Слой 1: память проекта & Техстек, стандарты, структура каталогов & Записывается в CLAUDE.md и подобные файлы, загружается автоматически, повторять не нужно \tabularnewline
Слой 2: ссылки на файлы & Файлы кода, затронутые текущей задачей & Ссылайтесь точно через символ @, не вставляйте файл целиком в диалог \tabularnewline
Слой 3: контекст диалога & История диалога по текущей задаче & Завершив задачу, открывайте новый диалог — контекст не должен разбухать \tabularnewline
Слой 4: текущая инструкция & Конкретные требования текущего шага & Ясно, кратко, по делу — без длинных описаний \tabularnewline
\hline
\end{tabularx}
\end{center}

\begin{center}\small\color{gray}Послойная стратегия управления окном контекста\end{center}

\booksection{Практические приёмы инженерии контекста}

Инженерия контекста — не разовая настройка, а процесс постоянной оптимизации. Вот несколько приёмов, многократно проверенных практикой:

\begin{enumerate}
  \item Обновляйте файл памяти вместе с эволюцией проекта. Появилась новая технология или изменилась архитектура — обновите файл памяти сразу, а не объясняйте «на ходу» в диалоге.
  \item Пользуйтесь ссылками @ вместо вставки целого текста. Большинство ИИ-инструментов поддерживают синтаксис @file или @folder: содержимое подставляется в контекст ИИ в фоне — эффективнее, чем вставлять вручную, и не засоряет диалог.
  \item Держите диалоги «короткими и ёмкими». Один диалог — одна задача или одна функциональная точка. Закончили — открывайте новый диалог для следующей. Длинные диалоги не только съедают окно контекста, но и «рассеивают внимание» ИИ.
  \item Давайте ИИ «увидеть» результат, а не «услышать описание». Наткнулись на баг — прикладывайте лог ошибки через ссылку @, а не пересказывайте своими словами «там вроде какая-то ошибка». ИИ полезнее видеть исходную информацию, чем ваш пересказ.
  \item Регулярно чистите и архивируйте. Завершённые диалоги и временные файлы занимают рабочее пространство. Периодическая уборка лишнего и порядок в каталоге проекта помогают ИИ эффективнее понимать структуру проекта.
\end{enumerate}

\begin{tcbpractice}{Практика: напишите для своего проекта 70-строчный файл памяти}
Возьмите проект, которым занимаетесь (или собираетесь), и напишите CLAUDE.md: обзор проекта (2–3 предложения), техстек (перечислите), стандарты кода (3–5 пунктов), структура каталогов (нарисуйте дерево), соглашения разработки (2–3 пункта). Уложитесь в 70 строк — закончив, вы заметите, что и сами стали понимать проект заметно яснее.

\end{tcbpractice}

\begin{tcbsummary}{Краткие итоги главы}
\begin{itemize}
  \item Память ИИ «краткосрочна»: в каждом новом диалоге он знает о проекте ноль
  \item Файл памяти проекта (CLAUDE.md и др.) — «долговременная память» ИИ, загружается автоматически в начале каждого диалога
  \item Хороший файл памяти = обзор + техстек + стандарты + структура каталогов + соглашения; 70 строк вполне достаточно
  \item SDD — разработка по спецификациям: сначала спецификация, потом код; спецификация — «единственный источник истины», она стабильнее кода
  \item Контекст управляется четырьмя слоями: память проекта ссылки @файлы история диалога текущая инструкция
  \item Пять практических приёмов: синхронно обновлять, ссылаться через @, короткие диалоги, давать исходные данные, регулярно чистить
\end{itemize}

\end{tcbsummary}
